\chapter{Introdução}

O solo constitui um sistema trifásico complexo formado por partículas sólidas, água e ar. Sob cargas externas, a interação entre essas fases governa o comportamento mecânico e hidráulico do subsolo. Na engenharia civil, o controle dessas resposta é o requisito fundamental para a estabilidade de fundações, barragens, aterros e pavimentos. Este trabalho analisa de forma concisa e rigorosamente fundamentada as metodologias de compactação e a adensamento de acordo com as normas brasileiras \citeonline{nbr7182} e \citeonline{mb33336}, descrevendo seus ensaios e parâmetros de projeto.